\documentclass[11pt]{article}
\usepackage[margin=1in]{geometry}
\usepackage{amsmath, amssymb}
\usepackage{enumitem}
\usepackage{titlesec}
\usepackage{hyperref}
\usepackage{booktabs}
\usepackage{xcolor}
\usepackage{parskip}

\hypersetup{colorlinks=true, linkcolor=blue, urlcolor=blue}

\title{\textbf{EE627 --- Class Summary}\\[0.4em]
       \large February 20, 2026}
\author{Prof.\ Rensheng Wang}
\date{}

\begin{document}
\maketitle
\tableofcontents
\newpage

%------------------------------------------------------------
\section{Closing the Loop: From ARIMA to Model Evaluation}
%------------------------------------------------------------

The lecture opened by connecting all prior time series topics into a single
workflow. The central question:

\textit{``You have two ARIMA models. Both were fit to the same data with
different $(p,d,q)$ parameters. Which is better?''}

The answer is \textbf{not} to look at the coefficients. The answer is to
look at the \textbf{residuals}.

\subsection{What Are Residuals --- Again}
For an AR(1) model $X_t = a_0 + a_1 X_{t-1} + \varepsilon_t$, after
estimation:
\[
\text{residual}_t = X_t - \left(\hat{a}_0 + \hat{a}_1 X_{t-1}\right)
\]
This is the portion of today's observation that the deterministic model
\emph{could not explain}. The model captures structure; the residual is
what is left over.

\subsection{The Residual as the Judge}
\textbf{Key insight:} if the model is perfect, the residuals must be
\textbf{pure white noise} --- IID, zero mean, constant variance, zero
autocorrelation for all lags.

\textit{Why?} The model's job is to extract every deterministic pattern
from the data. Noise, by definition, has no structure to extract. So if
any structure remains in the residuals, the model is inadequate --- it left
something exploitable behind.

\textit{Analogy from lecture:} squeezing an orange for juice. The model
must squeeze as hard as possible to extract every drop of signal. If the
residual still contains ``juice'' (structure), you have not squeezed hard
enough. The ideal residual is a dry peel --- pure noise, nothing more.

\subsection{How to Test Residuals}
A residual sequence is white noise if and only if:
\[
\text{ACF}(l) = 0 \quad\text{and}\quad \text{PACF}(l) = 0
\quad\text{for all } l > 0
\]
\textit{Procedure:}
\begin{enumerate}
    \item Fit ARIMA$(p,d,q)$.
    \item Extract \texttt{\$residuals}.
    \item Plot the ACF and PACF of the residuals.
    \item If all values lie within the confidence band (near zero) $\Rightarrow$
      adequate model.
    \item If spikes appear $\Rightarrow$ the model missed structure; revise
      $(p,d,q)$ and refit.
\end{enumerate}

\subsection{AIC for Choosing Among Adequate Models}
When \emph{multiple} models pass the residual test, use the
\textbf{Akaike Information Criterion (AIC)} to choose:
\[
\text{AIC} = 2k - 2\ln(\hat{L})
\]
where $k$ = number of estimated parameters and $\hat{L}$ = maximized
likelihood. \textbf{Smaller AIC = better model.}

AIC penalizes complexity --- a model that fits well with fewer parameters
is preferred. This prevents selecting an over-parameterized model that
memorizes noise rather than capturing genuine structure.

\textbf{Practical workflow:}
\begin{enumerate}
    \item Try multiple $(p,d,q)$ combinations; check residuals for each.
    \item From candidates that pass residual diagnostics, pick the one with
      the lowest AIC.
\end{enumerate}


%------------------------------------------------------------
\section{Reading ACF and PACF: The Diagnostic Table}
%------------------------------------------------------------

A major chunk of the lecture clarified how to identify the right model
family (AR vs.\ MA vs.\ ARMA) from plots. This is the key interview skill.

\subsection{Diagnostic Summary Table}

\begin{center}
\begin{tabular}{@{}llll@{}}
\toprule
\textbf{Model} & \textbf{ACF} & \textbf{PACF} & \textbf{Decision rule} \\
\midrule
AR($p$) & Gradual decay to zero & \textbf{Cliff at lag $p$} & Cliff in PACF \\
MA($q$) & \textbf{Cliff at lag $q$} & Gradual decay to zero & Cliff in ACF \\
ARMA($p,q$) & Gradual decay after lag $q$ & Gradual decay after lag $p$ & Both decay \\
White noise & All zeros & All zeros & No structure \\
\bottomrule
\end{tabular}
\end{center}

\textit{Where does ``cliff'' come from?} For an AR($p$) model, the PACF
value at lag $l$ is the last-added coefficient $\hat{a}_{l,l}$ from the
AR($l$) fit. For $l > p$, the ground truth says those lags do not exist,
so $\hat{a}_{l,l} \to 0$. The drop from nonzero to zero at lag $p$ is
the cliff.

\textit{Equivalently for MA:} MA($q$) = AR($\infty$) (from the AR--MA
duality), so its PACF never cuts off --- it decays gradually. But its
ACF cuts off sharply at $q$ because observations more than $q$ noise
terms apart share no common noise history.

\subsection{Identifying the Order from the Cliff}
\textbf{AR order $p$:} find the last lag in the PACF that lies
\emph{outside} the significance band. All subsequent lags should be inside.

\textbf{MA order $q$:} find the last lag in the ACF that lies
\emph{outside} the significance band. All subsequent lags should be inside.

The significance band is $\pm\, 2/\sqrt{T}$ (95\% level). Any PACF or
ACF value outside this band is statistically distinguishable from zero.

\textit{Key judgment:} if you observe 20 lags and lags 1 and 2 are outside
the band but lags 3--20 are inside, treat this as a ``cliff at lag 2'' and
specify order 2 --- even if a single later lag pokes slightly outside by
chance.

\subsection{Worked Example from Lecture (Homework Data)}
The professor showed a dataset whose ACF and PACF looked as follows:
\begin{itemize}
    \item \textbf{ACF:} lags 1 and 2 outside the band; lags 3 onward inside.
      $\Rightarrow$ cliff at lag 2 in the ACF.
    \item \textbf{PACF:} gradual decaying values, no clear cliff.
\end{itemize}

From the table: ACF cliff + PACF decay = \textbf{MA(2)} model.

\textit{Instructor reasoning:} ``If I see zeros from lag 3 onward, I trust
that the further lags will also be zero, because as the lag grows, the
probability of two time points sharing any informational link decreases.''


%------------------------------------------------------------
\section{Worked Example: Interest Rate Time Series}
%------------------------------------------------------------

The second example used real Yahoo Finance data (interest rate series).

\subsection{Step 1: Check Stationarity}
Plot the raw series. Observation: the mean is clearly changing over time
(trending). \textbf{Verdict: not stationary.} Cannot apply ACF/PACF to the
raw data.

\subsection{Step 2: First Differencing}
Compute $Y_t = X_t - X_{t-1}$.
\begin{itemize}
    \item The new series $Y_t$ has a mean that appears approximately
      constant (near zero).
    \item The variance changes slightly in different windows, but not
      dramatically --- close enough to treat as stationary for this exercise.
\end{itemize}

\subsection{Step 3: Fit ARIMA(3,1,2)}
Based on inspection of the ACF/PACF of $Y_t$, the professor tried
ARIMA(3,1,2): AR order 3, one difference, MA order 2.

\subsection{Step 4: Check Residuals}
After fitting, plot ACF and PACF of \texttt{\$residuals}:
\begin{itemize}
    \item ACF: lag 0 = 1 (expected); lags 1 onward approximately within
      the blue significance bands.
    \item PACF: mostly within bounds, with minor deviations at the
      far end.
\end{itemize}
\textbf{Verdict:} residuals are close to white noise. ARIMA(3,1,2) is an
adequate model for this series. Further tuning could improve it, but the
baseline is already good.

\textit{Key message:} in a real interview or job setting, the boss does not
look at whether you called the right function. The boss plots your residuals.
If the residuals are not noise, the model is rejected --- regardless of how
elegant the formula looks.


%------------------------------------------------------------
\section{Recommender Systems: Business Motivation}
%------------------------------------------------------------

\subsection{Why Recommenders Matter}
\begin{itemize}
    \item \textbf{Unavoidable:} every click, every scroll, every search
      result is mediated by a recommender. Google, YouTube, Amazon, Netflix,
      TikTok, Instagram --- all run on recommenders.
    \item \textbf{Two sides of the same product:} recommenders and search
      engines are mirror images. ``Flip a recommender around and you get a
      search engine.'' Both match a user to relevant content.
    \item \textbf{Instant monetization:} most industries require a long
      production cycle before revenue arrives (plant soybeans, wait a year).
      Recommenders generate cash in \textit{milliseconds} --- the moment
      you click, money flows. This is why companies like Meta and Google make
      their money through recommenders, not software licenses.
    \item \textbf{Compensation premium:} engineers working on recommender
      and search systems earn dramatically more than typical software
      engineers, because a 0.001\% improvement at billion-user scale
      translates to enormous profit.
\end{itemize}

\subsection{Recommenders Shape Reality}
The professor made a point worth recording: recommenders do not only
\emph{reflect} user preferences --- they actively \emph{shape} them.
\begin{itemize}
    \item A recommender that always returns exactly what you searched for
      creates a closed feedback loop (``filter bubble''): you watch football,
      you get football, you watch more football --- you never discover
      anything new.
    \item Good recommenders deliberately add \textbf{diversity} --- a small
      fraction of results that go slightly beyond what you searched for ---
      to develop new interests and prevent the loop from closing.
    \item This is not an accident or an error; it is intentional psychological
      design to expand the user's engagement surface.
\end{itemize}


%------------------------------------------------------------
\section{The Four-Stage Recommender Architecture}
%------------------------------------------------------------

Modern recommender systems at scale (YouTube, Amazon, Spotify) process
\textbf{billions of items} to deliver a ranked list of 10--20 results in
milliseconds. A single-stage model cannot do this. The standard architecture
uses four cascading stages, each narrowing the candidate pool.

\begin{center}
\begin{tabular}{@{}lll@{}}
\toprule
\textbf{Stage} & \textbf{Candidate pool} & \textbf{Primary goal} \\
\midrule
1. Retrieval / Candidate Generation & Billions $\to$ Thousands & Speed; broad relevance \\
2. Rough Ranking (Pre-filter)        & Thousands $\to$ Hundreds & Rule-based filtering \\
3. Fine Ranking                      & Hundreds $\to$ Tens      & Precision scoring \\
4. Final Re-ranking                  & Tens $\to$ Top 10        & Business rules + diversity \\
\bottomrule
\end{tabular}
\end{center}

\subsection{Offline vs.\ Online Pipeline}
The architecture runs as two separate pipelines:

\begin{itemize}
    \item \textbf{Offline pipeline:} pre-computes and stores item embeddings
      for all items in a \textbf{Vector Database (Vector DB)}. This runs
      continuously in the background, not in real time.
    \item \textbf{Online pipeline:} when a user makes a request, computes
      the customer embedding in real time (using the trained Two-Tower
      model), queries the Vector DB for nearest neighbors, then passes
      those candidates through Stages 2--4.
\end{itemize}

The offline/online split is what makes billion-item retrieval fast: the
expensive pre-computation is done once and stored; only the customer query
runs live.

\subsection{Stage 1: Retrieval / Candidate Generation}
\textbf{Goal:} reduce from $>$100M--billions of items to $\sim$100--1000
relevant candidates per user.

\textbf{Key approach --- embedding-based Two-Tower retrieval:}
\begin{itemize}
    \item Every user and every item is mapped to a \textbf{dense embedding
      vector} in the same latent space.
    \item These vectors encode semantic meaning numerically --- similar
      items cluster together in the space.
    \item Retrieval = approximate nearest-neighbor (ANN) search: find the
      item vectors closest to the customer query vector.
    \item Similarities expressed initially in coarse features (customer ID,
      item ID); fine attributes are added in later stages.
\end{itemize}

\textbf{Retrieval methods:}
\begin{itemize}
    \item \textbf{Content-based:} match user/item embeddings via dot product
      or cosine similarity.
    \item \textbf{Collaborative filtering:} user-item interaction matrix
      (matrix factorization, ANN search).
    \item \textbf{Graph-based:} random walks on the user-item interaction
      graph to find structurally proximate candidates.
    \item \textbf{Heuristic rules:} popularity, trending, recency boosts
      (effective for cold-start and smaller datasets).
\end{itemize}

\textbf{Output:} $\sim$100--1000 candidate items per user.

\subsection{Stage 2: Rough Ranking (Filtering / Preliminary Ranking)}
\textbf{Goal:} filter out irrelevant or ineligible candidates cheaply,
before the expensive Stage 3 scoring. Output: $\sim$100--300 items.

\textbf{Key tasks:}
\begin{itemize}
    \item Remove items the user has already seen or purchased.
    \item Remove unavailable, policy-violating (NSFW), or out-of-region
      items.
    \item Apply basic scoring for click probability or engagement using
      \emph{shallow} ML models (logistic regression, gradient boosting).
    \item Filter out irrelevant categories based on context.
\end{itemize}

\textbf{Bloom filter:} seen/purchased item removal is often implemented
with a \textbf{Bloom filter} --- a space-efficient probabilistic data
structure that answers ``has this item been seen before?'' in $O(1)$
time with very low memory, at the cost of a small false-positive rate.

\subsection{Stage 3: Fine Ranking}
\textbf{Goal:} assign precise relevance scores to each
\texttt{(item, customer)} pair and sort by score. Output: top $\sim$20--50
items.

\textbf{Why this stage is expensive:} unlike Stage 1 (fast approximate
similarity), Stage 3 uses rich features for every candidate:
\begin{itemize}
    \item Item features: title, tags, description, category, rating.
    \item Customer features: demographic info, historical behavior,
      contextual signals (device, time, location).
    \item All features fetched on demand from a \textbf{feature store}
      (a low-latency key-value store for pre-computed features).
\end{itemize}

\textbf{Models used:}
\begin{itemize}
    \item XGBoost or CatBoost (for tabular features)
    \item Deep neural networks (DNNs), e.g., YouTube's ranking network
    \item Multi-objective optimization targets: CTR, CVR, watch time,
      revenue, diversity
    \item LLMs (emerging use case)
\end{itemize}

After scoring, items are ordered by their ranking score, optionally cut to
$xK$ items before Stage 4. Having the most personalized item at rank 1 is
critical --- users rarely scroll past the first few results.

\subsection{Stage 4: Re-ranking (Post-ranking)}
\textbf{Goal:} optimize the final ordered list for multiple objectives
beyond pure relevance. Output: final list shown to the user.

\textbf{Key adjustments:}
\begin{itemize}
    \item \textbf{Diversity / novelty / serendipity:} deliberately inject
      items slightly outside the user's usual interest to prevent filter
      bubbles and develop new engagement. Algorithms: MMR (Maximal Marginal
      Relevance), xQuAD, fairness-aware ranking.
    \item \textbf{Sponsored content:} blend in advertisements according to
      business contracts and content mix policies.
    \item \textbf{Multi-source blending:} merge results from multiple
      recommendation signals (``Because you watched\ldots'' +
      ``Trending now'' + ``Friends liked\ldots'').
    \item \textbf{Exploration:} bandit or reinforcement learning to
      adaptively reorder and discover what works over time.
    \item \textbf{Multi-objective trade-offs:} balance engagement vs.\
      satisfaction vs.\ fairness vs.\ revenue.
\end{itemize}

\textit{Key insight:} final re-ranking is \emph{not} purely a ranking
problem --- it is a multi-objective optimization that reflects business
strategy, psychology, and fairness constraints. This is where engineering
judgment --- not just model output --- determines the product experience.

\subsection{Why Only the Top 10 Matter}
Users do not scroll beyond the first 10--20 results. A recommendation at
position 200 has near-zero click probability. Each stage's primary job is
to preserve the highest-quality candidates as the pool shrinks, so the
final top-10 list is maximally relevant.


%------------------------------------------------------------
\section{The Two-Tower Model: Architecture and Training}
%------------------------------------------------------------

The two-tower model is the standard embedding architecture for Stage 1
retrieval. Understanding it is important because the same structure appears
throughout deep learning (attention in transformers, contrastive learning,
semantic search).

\subsection{Architecture}
Two neural networks (``towers'') run in parallel, each producing a
dense embedding vector:

\begin{center}
\begin{tabular}{@{}lll@{}}
\toprule
\textbf{Tower} & \textbf{Inputs} & \textbf{Output} \\
\midrule
Customer Query Encoder & Demographics, historical behavior, contextual features & $\mathbf{u} \in \mathbb{R}^d$ \\
Item Candidates Encoder & Tags, description, rating, category & $\mathbf{v} \in \mathbb{R}^d$ \\
\bottomrule
\end{tabular}
\end{center}

Both towers project into the \textbf{same $d$-dimensional embedding space}.
This is the key constraint --- it is what allows direct comparison.

\subsection{Similarity and Loss}
The predicted relevance of item $j$ for customer $i$ is the dot product:
\[
\text{score}(i, j) = \mathbf{u}_i \cdot \mathbf{v}_j
\]
\begin{itemize}
    \item Vectors nearly parallel $\Rightarrow$ large dot product $\Rightarrow$
      high relevance.
    \item Vectors orthogonal $\Rightarrow$ dot product near zero $\Rightarrow$
      no match.
\end{itemize}

\textbf{Training signal:} the interaction matrix provides labels.
\begin{itemize}
    \item Positive pair (customer bought/watched/clicked item): target
      score = 1.
    \item Negative pair (no interaction): target score = 0.
\end{itemize}
The loss function minimizes the difference between the dot product scores
and these binary labels, forcing matched pairs to be close in embedding
space and non-matched pairs to be far apart.

\subsection{Interaction Matrix}
The training data is the \textbf{user-item interaction matrix}:
rows = users, columns = items, cells = feedback signal.
\begin{itemize}
    \item \textbf{Explicit feedback:} star ratings (1--5), thumbs up/down.
    \item \textbf{Implicit feedback:} watch time, click, purchase, share,
      dwell time. More common in practice because explicit ratings are
      sparse.
\end{itemize}

\subsection{Why Dot Product / Why Latent Space}
The computer has no semantic understanding of ``swimming goggle.''
In original item-ID space, item 10033 and item 10034 are just two numbers
--- there is no notion of similarity. In the learned latent space, similar
items cluster together numerically. The dot product measures how aligned
two vectors are --- a purely numerical operation the computer executes
instantly. This is why embedding + ANN search can retrieve relevant items
from a billion-item corpus in milliseconds.

The same inner-product similarity is used in transformers (attention
= $\mathbf{Q}\mathbf{K}^T$), contrastive learning (CLIP, SimCLR), and
semantic search. The two-tower model is a foundational building block
across modern AI.

\end{document}
